%% Figure block, shared by ../paper/main.tex and ../submission/*.tex so the caption has
%% exactly one source.
\begin{figure}[tbp]
\centering
\figorbox{fig7_lrv_rate.pdf}
\caption{\textbf{The rate claim, measured, on the object an interval actually uses.}
Both axes logarithmic; error bars are $\pm1.96$ Monte Carlo standard errors; fitted slopes in
the legends. \textbf{(a)} Relative operator-norm error in estimating the whole sandwich
$\Sigma=H^{-1}SH^{-1}$. Ours is $\Hb^{-1}\Sft\Hb^{-1}$ at $b\asymp\log N$; the comparator is
overlapping batch means on the \emph{iterate} path at the block length $b_n\asymp T^{3/4}$ that
analysis requires (\Cref{app:hyper} gives the form), which uses no gradients and no curvature
estimate. Same target, so the comparison is meaningful; but read the levels, not the slopes. The
iterate-path curve is flat here, and its published $N^{-1/8}$ rate is too slow for our range to
resolve, so its fitted slope says nothing about it. What the panel shows is that over the sample
sizes we can reach, one estimator is accurate and the other is not. \textbf{(b)} Error in estimating $S$ alone, each estimator at its own optimal block
length: the two-scale form at $b\asymp\log N$ against the $-1/2$ of \Cref{thm:lrv}(iv), plain
batch means at $b\asymp N^{1/3}$ against $-1/3$.}
\label{fig:lrvrate}
\end{figure}
